\chapter{Kết luận và hướng nghiên cứu}
\index{encoding!lựa chọn}
\index{bằng chứng thực nghiệm}
\index{giới hạn phương pháp}
\index{chứng nhận tối ưu}

\chapterlead{Chương kết nối các kết quả kỹ thuật thành một luận điểm chung:
một phép biểu diễn SAT tốt là một giao diện có ngữ nghĩa, được chọn theo cấu trúc
bài toán và được đánh giá bằng một chuỗi bằng chứng có thể kiểm tra. Từ đó,
chương tổng hợp miền hiệu quả của từng họ phương pháp và xây dựng chương trình
nghiên cứu tiếp theo.}

\section{Những kết luận bền vững}

Các họ bộ biểu diễn tạo thành một biên Pareto (Pareto frontier) thay vì một thứ hạng duy nhất. Quy
trình lựa chọn có thể dùng lại gồm bốn lớp.

\paragraph{Lớp ngữ nghĩa.}
Trước khi đếm biến và mệnh đề, phải viết đặc tả cho biến chính, biến phụ, bộ
giải mã và hàm mục tiêu. Tính đúng, tính đầy đủ theo phép chiếu và khả năng dùng
literal phụ ở mô-đun khác là ba yêu cầu khác nhau. Một biến Tseitin chỉ được
xem như ngưỡng, biến làm chứng hay trạng thái khi các mệnh đề thật sự bảo đảm ngữ
nghĩa đó.

\paragraph{Lớp cấu trúc.}
Phép biểu diễn nên khớp với hình học của ràng buộc. Cửa sổ chồng lấn gợi ý trạng
thái chia sẻ; miền hữu hạn có thứ tự gợi ý literal thứ tự; tổng Boolean gợi
ý bộ đếm, cây tổng, mạng đếm hoặc BDD; phép tuyển hình học gợi ý biến quan hệ.
Đối xứng nên được loại trước hết ở cách chọn biến, sau đó mới bằng ràng buộc có
phép ánh xạ bảo toàn nghiệm.

\paragraph{Lớp giải.}
Đánh giá ở lớp giải kết hợp số mệnh đề với lan truyền đơn vị, đường truyền qua
biến phụ, khả năng giữ mệnh đề đã học khi thay cận, chi phí của bộ giải mã/bộ
kiểm tra và tương tác với tiền xử lý. Cùng một phép biểu diễn khả thi có thể cư xử
khác nhau dưới SAT độc lập, SAT tăng dần, MaxSAT hay một kiến trúc kết hợp.

\paragraph{Lớp bằng chứng.}
Phát biểu về tính đúng cần chứng minh hoặc kiểm tra vét cạn trên miền nhỏ; phát
biểu về kích thước cần công thức và bộ đếm từ gói tái lập; phát biểu về hiệu
năng cần quy trình, đối chứng, giới hạn thời gian và dữ liệu từng lần chạy.
Tối ưu chỉ được chứng nhận khi cận trên từ một nghiệm khả thi gặp cận dưới hợp
lệ, thường là một kết quả
\UNSAT ở cận kế tiếp.

\section{Tổng hợp các nghiên cứu tình huống}

Bảng~\ref{tab:reported-evidence-map} tổng hợp thiết kế đánh giá, kết quả và cơ
chế biểu diễn nổi bật của các nghiên cứu tình huống.

\begin{longtable}{@{}
>{\raggedright\arraybackslash}p{.19\textwidth}
>{\raggedright\arraybackslash}p{.23\textwidth}
>{\raggedright\arraybackslash}p{.39\textwidth}@{}}
\caption{Tổng hợp các nghiên cứu tình huống.}
\label{tab:reported-evidence-map}\\
\toprule
Mô-đun/bài toán & Thiết kế đánh giá & Kết quả chính\\
\midrule
\endfirsthead
\toprule
Mô-đun/bài toán & Thiết kế đánh giá & Kết quả chính\\
\midrule
\endhead
SCL/SCAMO và Anti-Bandwidth &
Phân tích kích thước SCAMO; 24 đồ thị kiểm chuẩn &
Thang chia sẻ giảm phần bộ đếm lặp trên AMO bậc thang và cạnh tranh trong
cấu hình Anti-Bandwidth đã thử; độ chồng lấn khoảng là đặc trưng quyết định
\parencite{truong2025scamo}.\\

NSC/Xếp nhóm chơi golf &
66 bài toán; NSC, ADDER, BDD, CARD, BINOMIAL &
NSC giải \(56/66\); ADDER, BDD và CARD cùng \(54/66\), BINOMIAL \(32/66\).
Kết quả định vị NSC như một lựa chọn hiệu quả cho ràng buộc đúng \(K\) trong mô hình
bài toán xếp nhóm chơi golf của nguồn \parencite{kieu2026golfer}.\\

ALSCE/Lập lịch điều dưỡng &
28 bài toán, đến 120 điều dưỡng và 24 tuần &
ALSCE giải toàn bộ tập; kết quả hỗ trợ việc chia sẻ thanh ghi cho các cửa sổ
chuỗi trong mô hình này \parencite{truong2025nurse}.\\

SAT/Đóng gói dải &
41 bài toán; SAT thường, SAT tăng dần và CP &
Số giá trị tối ưu được chứng nhận lần lượt là \(26\), \(21\) và \(28\). SAT
tăng dần phát huy lợi thế khi thay cận giữ được phần lớn miền và mệnh đề đã học
\parencite{kieu2025strip}.\\

Gán nhãn radio &
146 bài toán từ chín họ đồ thị; các khung SAT và ILP &
Công trình báo 109 giá trị tối ưu được chứng nhận, 38 số radio mới, và khớp
hoặc cải thiện BKS trên 130 bài toán; SAT và ILP có vai trò bổ sung theo cấu
trúc đồ thị
\parencite{thanh2026radio}.\\

Gán nhãn an toàn \(k\) &
16 đồ thị Erdős--Rényi; FE, RE và DE &
FE và RE tìm nghiệm cho \(16/16\), chứng nhận 13 giá trị tối ưu; DE tìm nghiệm
cho 15 và chứng nhận 8. Biên thứ tự phù hợp hơn biểu diễn trực tiếp trong tập thử
này \parencite{thanh2026ksafe}.\\

POSE/MO-FAP &
35 bài toán CALMA/DUTtest hoặc được mô hình hóa lại &
POSE kết hợp bộ đếm tăng dần quyết định toàn bộ tập thử; kết quả hỗ trợ
kiến trúc lọc miền, nén khoảng cấm và đếm số tần số khác nhau
\parencite{dao2026mofap}.\\
\bottomrule
\end{longtable}

\section{Miền hiệu quả và điểm chuyển thiết kế}

\paragraph{Bộ đếm chia sẻ và Ladder AMK.}
Bộ đếm chia sẻ được ưu tiên khi các cửa sổ chồng lấn mạnh, cận \(k\) nhỏ và
mỗi cửa sổ cắt ít khối. Khi số biên tăng, bộ đếm ghép thay bộ nối hai ngôi;
với các tập gần rời nhau, một phép biểu diễn cục bộ thường gọn hơn.

\paragraph{NSC.}
NSC có kích thước \(O(nk)\) và tiết kiệm vùng trạng thái tam giác bất khả. Cấu trúc này
phù hợp nhất khi vùng này chiếm tỷ lệ đáng kể và bài toán dùng một cận cố định.
Cây tổng/mạng đếm phù hợp với nhiều cận tăng dần; khi \(k\) gần \(n\), đếm
literal bù chuyển bài toán về ngưỡng \(n-k\) nhỏ hơn.

\paragraph{Phá vỡ đối xứng.}
Mỗi ràng buộc phá vỡ đối xứng được suy từ một nhóm tự đẳng cấu bảo toàn miền, ràng
buộc và hàm mục tiêu. Sức chứa, miền quay và cách hàm mục tiêu đọc chỉ số quyết
định những hoán vị nào thuộc nhóm này. Mức phá vỡ đối xứng được chọn cùng độ dài
mệnh đề và chiến lược phân nhánh.

\paragraph{Biểu diễn thời gian và đóng gói.}
SAT theo mốc thời gian phù hợp với chân trời vừa phải và miền bắt đầu đã được
lọc. Tọa độ thứ tự nén tốt các so sánh trên miền nguyên. DDD mở rộng lựa chọn
cho chân trời lớn khi có tập mốc đầy đủ hữu hạn, cận dưới hợp lệ và bước tinh
chỉnh tăng nghiêm ngặt.

\paragraph{Khoảng cấm và biên thứ tự.}
Hai literal biên nén hiệu quả một khoảng cấm trên miền có thứ tự và đặc tả thứ
tự hai chiều. Miền rời rạc dùng ánh xạ phần tử trước/sau; mỗi khoảng rời nhau
bổ sung một cặp biên. Nhất quán cung tạo nhiều giá trị nhất trước SAT khi miền
lớn hoặc đồ thị ràng buộc đủ dày.

\section{Hướng nghiên cứu tiếp theo}

\begin{enumerate}
  \item \textbf{Lựa chọn bộ biểu diễn theo bài toán.} Xây bộ chọn có thể giải thích
  quyết định bằng các đặc trưng như \(n,k\), độ chồng lấn, độ thưa của miền và
  số cận sẽ truy vấn; đánh giá bộ chọn so với bộ giải tốt nhất ảo.
  \item \textbf{Đặc tả và kết hợp.} Phát triển thư viện mô-đun mà mỗi literal
  đầu ra có ngữ nghĩa máy đọc được, cho phép kiểm tự động việc ghép bộ đếm,
  liên kết biểu diễn, hàm mục tiêu và bộ giải mã.
  \item \textbf{Chứng nhận toàn trình.} Nối ghi chứng minh của SAT/MaxSAT với
  bộ kiểm tra miền và chứng nhận cận dưới để giá trị tối ưu không phụ thuộc vào niềm tin ở
  bộ sinh.
  \item \textbf{Trạng thái chia sẻ tổng quát.} Mô hình hóa việc chọn khối,
  cây ghép và ngưỡng cần hiện thực hóa (reification) như một bài toán tối ưu riêng; tìm cận dưới
  cho số trạng thái/bộ nối dưới các yêu cầu lan truyền khác nhau.
  \item \textbf{Phương pháp kết hợp có phân công rõ.} Dùng CP/MIP để lọc miền
  hoặc tạo cận, SAT/MaxSAT để xử lý phép tuyển, nhưng đánh giá bằng thí nghiệm
  phân tích thành phần (ablation) để biết lợi ích đến từ mô-đun nào.
  \item \textbf{Gói tái lập bền vững.} Dùng vùng chứa phần mềm, bản kê dữ liệu,
  mã băm, giấy phép và quy trình tái lập; tách kết luận đã công bố khỏi kết quả
  thăm dò.
\end{enumerate}

\section{Kết luận cuối}

Giá trị của phép biểu diễn SAT nằm ở khả năng chọn đúng trạng thái trung gian để
CDCL nhìn thấy cấu trúc của bài toán. Bộ đếm chia sẻ, thanh ghi tam giác,
biên thứ tự, biến quan hệ và rời rạc hóa động là các biểu hiện khác nhau của
cùng một nguyên tắc: chỉ hiện thực hóa (reification) thông tin vừa đủ, đặt đặc tả rõ cho thông tin
đó, rồi kiểm bằng cả chứng minh lẫn gói tái lập. Ba bước này chuyển một
nghiên cứu tình huống thành một phương pháp có thể kiểm chứng, so sánh và dùng
lại.
